\documentclass[answers]{exam}
\usepackage{../../../mypackages}
\usepackage{../../../macros}

\title{Interrogation calcul mental}
\author{N. Bancel}
\date{2 Décembre 2025}

% Mise en forme des solutions en bleu
\SolutionEmphasis{\color{blue}}
\renewcommand{\solutiontitle}{\noindent}

\begin{document}

\textbf{Collège Lycée Suger}
\hfill
\textbf{Mathématiques} \

\textbf{Année 2025-2026}
\hfill
\textbf{Classe de 6ème} \par

{\let\newpage\relax\maketitle}

\begin{center}
\textbf{\textcolor{red}{Correction de l'interrogation N°4 de calcul mental}} \\
Les solutions en \textcolor{blue}{bleu} expliquent une méthode possible pour trouver le résultat.
\end{center}

\begin{questions}

\question $40 + 5.8 + 0.2 + 9 =$
\begin{solution}
On cherche des regroupements astucieux pour obtenir des nombres ronds (multiples de 10 si possible). On utilise la commutativité et l'associativité de l'addition pour permuter et regrouper les termes.
\begin{align}
40 + 5.8 + 0.2 + 9
&= 40 + 9 + 5.8 + 0.2 &&\text{par commutativité de l'addition} \\
&= (40 + 9) + (5.8 + 0.2) &&\text{par associativité de l'addition} \\
&= 49 + 6.0 &&\text{car } 5.8 + 0.2 = 6.0 \\
&= 55.0 &&\text{car } 49 + 6 = 55
\end{align}
Donc $40 + 5.8 + 0.2 + 9 = 55$.
\end{solution}

\question $76 \div 4 =$
\begin{solution}
On utilise l'astuce : « diviser par 4, c'est diviser deux fois par 2 ». On effectue les étapes et on écrit chaque division.
\begin{align*}
76 \div 4
&= 76 \div (2 \times 2) &&\text{on remplace }4\text{ par }2\times2 \\
&= (76 \div 2) \div 2 &&\text{on divise d'abord par }2 \\
&= 38 \div 2 &&\text{car }76 \div 2 = 38 \text{ (et que } 38 \times 2 = 76\text{)} \\
&= 19. &&\text{car }38 \div 2 = 19 \text{ (et que } 19 \times 2 = 38\text{)}
\end{align*}
Vérification rapide : $19 \times 4 = (19 \times 2) \times 2 = 38 \times 2 = 76$.
Donc $76 \div 4 = 19$.
\end{solution}

\question $(12 - 3) \times 2 + 4 =$
\begin{solution}
On applique la règle de priorité : on commence par les parenthèses, puis la multiplication avant l'addition.
\begin{align*}
(12 - 3) \times 2 + 4
&= 9 \times 2 + 4 &&\text{car }12 - 3 = 9 \text{ (on commence par les parenthèses)}  \\
&= 18 + 4 &&\text{car }9 \times 2 = 18 \text{ (priorité multiplication)} \\
&= 22. &&\text{addition finale}
\end{align*}
Donc $(12 - 3) \times 2 + 4 = 22$.
\end{solution}

\question $84162 \div 100 =$
\begin{solution}
Diviser par $100$ revient à diviser par $10$ puis encore par $10$ : la virgule se décale de deux rangs vers la gauche. On raisonne étape par étape et on vérifie ensuite par multiplication.
\begin{align*}
84162 \div 100
&= 84162 \div 10 \div 10 &&\text{car }100 = 10 \times 10 \\
&= 8416.2 \div 10 &&\text{la virgule recule d'un rang} \\
&= 841.62 &&\text{la virgule recule d'un deuxième rang}
\end{align*}
Vérification par multiplication : \
\begin{align*}
841.62 \times 100 &= 84162. &&\text{la virgule avance de deux rangs}
\end{align*}
Donc $84,162 \div 100 = 841.62$.

\textbf{\textcolor{red}{On divise 84162 en 1000 parts égales : chaque part vaut donc beaucoup moins que 84152. On ne peut pas tomber sur un nombre plus grand}} \\
\end{solution}

\question $117 \times 4 =$
\begin{solution}
On utilise l'astuce « multiplier par 4, c'est multiplier deux fois par 2 » et l'associativité de la multiplication.
\begin{align*}
117 \times 4
&= 117 \times (2 \times 2) &&\text{on remplace }4\text{ par }2\times2 \\
&= (117 \times 2) \times 2 &&\text{par associativité} \\
&= 234 \times 2 &&\text{car }117 \times 2 = 234 \\
&= 468. &&\text{car }234 \times 2 = 468
\end{align*}
Donc $117 \times 4 = 468$.
\end{solution}

\question $88.6 \times 1000 =$
\begin{solution}
Multiplier par $1,000$ équivaut à multiplier par $10$ trois fois : la virgule se décale de trois rangs vers la droite.
\begin{align*}
88.6 \times 1,000
&= 88.6 \times 10 \times 10 \times 10 &&\text{car }1,000=10\times10\times10 \\
&= 886 \times 10 \times 10 &&\text{la virgule avance d'un rang} \\
&= 8,860 \times 10 &&\text{la virgule avance encore} \\
&= 88,600. &&\text{la virgule avance d'un troisième rang}
\end{align*}
Donc $88.6 \times 1000 = 88,600$.
\end{solution}

\question $12 \div 1000 =$
\begin{solution}
Diviser par $1,000$ revient à diviser par $10$ trois fois : la virgule se décale de trois rangs vers la gauche. On écrit d'abord $12$ comme $12.0$ pour voir la virgule.
\begin{align*}
12 \div 1,000
&= 12.0 \div 10 \div 10 \div 10 &&\text{car }1,000=10\times10\times10 \\
&= 1.2 \div 10 \div 10 &&\text{la virgule recule d'un rang} \\
&= 0.12 \div 10 &&\text{la virgule recule d'un deuxième rang} \\
&= 0.012. &&\text{la virgule recule d'un troisième rang}
\end{align*}
Vérification : $0.012 \times 1,000 = 12$.
Donc $12 \div 1000 = 0.012$.
\end{solution}

\question $4 + 2 \times (5 - 1) =$
\begin{solution}
Respect des priorités : d'abord la parenthèse, puis la multiplication, puis l'addition.
\begin{align*}
4 + 2 \times (5 - 1)
&= 4 + 2 \times 4 &&\text{car }5 - 1 = 4 \\
&= 4 + 8 &&\text{car }2 \times 4 = 8 \\
&= 12. &&\text{addition finale}
\end{align*}
Remarque : la multiplication est prioritaire sur l'addition, on ne fait pas $(4+2)\times4$.
Donc $4 + 2 \times (5 - 1) = 12$.
\end{solution}

\question $34.18 \times 1000 =$
\begin{solution}
Multiplier par $1,000$ décale la virgule de trois rangs vers la droite.
\begin{align*}
34.18 \times 1,000
&= 34.18 \times 10 \times 10 \times 10 &&\text{car }1,000=10\times10\times10 \\
&= 341.8 \times 10 \times 10 &&\text{la virgule avance d'un rang} \\
&= 3,418 \times 10 &&\text{la virgule avance encore} \\
&= 34,180. &&\text{la virgule avance d'un troisième rang}
\end{align*}
Donc $34.18 \times 1000 = 34,180$.
\end{solution}

\question $0.5 \times 44 \times 4 =$
\begin{solution}
On utilise la commutativité et l'associativité de la multiplication pour regrouper les facteurs de manière astucieuse. Ici on remarque que $0.5 \times 4 = 2$, ce qui rend le calcul plus simple.
\begin{align*}
0.5 \times 44 \times 4
&= (0.5 \times 4) \times 44 &&\text{par commutativité et associativité} \\
&= 2 \times 44 &&\text{car }0.5 \times 4 = 2 \\
&= 88. &&\text{car }2 \times 44 = 88
\end{align*}
Alternative possible : regrouper $44 \times 4 = 176$ puis multiplier par $0.5$ pour obtenir $88$ (même résultat).
Donc $0.5 \times 44 \times 4 = 88$.
\end{solution}

\end{questions}

\end{document}
